\chapter{泛型}
\label{ch:14}

如果一段算法只依赖少数操作，却要用于多种元素类型，可以把变化的类型写成参数。泛型保留静态检查，并在具体使用时结合类型、常量和约束确定合法实现。

本章依次区分类型参数、整数常量参数、关联类型和泛型接口，再把控制表达式放入泛型函数中。阅读时始终说明每个参数代表什么，以及它在实例化后约束了哪些字段、方法和结果。

\section{类型参数和常量参数}

先想一个固定容量缓冲区：需要变化的既有元素类型，也有槽位数量。给每种组合手写一个结构体会重复相同字段与方法，因此把变化部分放到类型参数表中。

Buffer 用 T 表示元素类型，用带 int64 类型的 N 表示常量维度。二者都要在具体使用时确定，却不是同一种参数；下面通过 data 的形状与 len 的返回值分别观察它们。

\par
\medskip
\begin{lstlisting}[style=chali]
struct Buffer<T, N: int64>
{
    T data[N];
    int64 len()
    {
        return N;
    }
};
void main()
{
    Buffer<int64, 4> buffer;
    buffer.data[0] = 7;
    println("{} {}", buffer.len(), buffer.data[0]);
}
\end{lstlisting}
\par
\medskip

\par
\medskip
\begin{lstlisting}[style=out]
4 7
\end{lstlisting}
\par
\medskip

\subsection*{代码详解}

Buffer 的 T 是元素类型参数，N 是 int64 常量参数。声明 data[N] 时，N 参与数组形状；它不是每个对象可随意修改的长度字段。len 返回这个类型参数给定的容量值。

\texttt{Buffer<\allowbreak{}int64,\allowbreak{} 4>\allowbreak{}} 把两者具体化为整数元素和四个槽位，写入下标 0 后输出 4 与 7。数组合法下标为 0 到 3，len 的返回值 4 不是可访问的最后下标。选择不同 N 会改变实例化后的类型与布局，不能把它当成一次普通运行期赋值。

\section{以约束表达能力}

泛型代码不能假定任意 T 都能相加、哈希或读取某个字段。把所需行为写进约束；关联类型表达容器等接口依赖的元素类型。这样调用者看到的是能力要求，而不是某个偶然使用的实现字段。

\par
\medskip
\begin{lstlisting}[style=chali]
interface Container
{
    type Item;
    Item get(int64 index);
};
struct Numbers: Container
{
    type Item = int64;
    Item get(int64 index)
    {
        return index + 10;
    }
};
C.Item first<C: Container>(C value)
{
    return value.get(0);
}
void main()
{
    Numbers numbers;
    println("{}", first<Numbers>(numbers));
}
\end{lstlisting}
\par
\medskip

\par
\medskip
\begin{lstlisting}[style=out]
10
\end{lstlisting}
\par
\medskip

\subsection*{代码详解}

Container 声明关联类型 Item，把元素类型留给实现者决定。Numbers 用 \texttt{type Item =\allowbreak{} int64} 补全这个约定，其 get 返回 index+10，因此索引 0 对应 10。

first 的泛型约束 \texttt{C:\allowbreak{} Container} 保证可以调用 get；返回类型 \texttt{C.\allowbreak{}Item} 则把结果与同一个实现的关联类型连接起来。main 显式指定 Numbers，编译器据此得到 int64 返回值。关联类型不是字符串形式的类型名，也不是可以在运行时任意切换的字段。

接口本身也能参数化。Source<int64> 约定读取整数，ScoreSource 提供具体实现，使用者通过接口调用。

\par
\medskip
\begin{lstlisting}[style=chali]
interface Source<T>
{
    T read();
};
struct ScoreSource : Source<int64>
{
    int64 read() { return 42; }
};
void printScore(Source<int64> source)
{
    println("{}", source.read());
}
void main()
{
    ScoreSource source;
    printScore(source);
}
\end{lstlisting}
\par
\medskip

\par
\medskip
\begin{lstlisting}[style=out]
42
\end{lstlisting}
\par
\medskip

\subsection*{代码详解}

Source<T> 把 read 的返回类型参数化。ScoreSource 实现的是具体的 Source<int64>，因此必须提供返回 int64 的 read；示例直接返回 42，未读取文件或输入。

printScore 接收 Source<int64> 接口值，main 传入符合这个接口的实现。函数体通过接口调用 read，再打印 42。这里固定的是接口的类型实参，不能据此把 Source<String> 当作同一参数类型传入，也不需要为每种具体实现复制 printScore 的业务逻辑。

\section{返回值与控制流}

返回类型仍然约束所有可达路径。泛型函数里的 if、match、循环与普通函数遵守相同的结果规则，不能因为类型尚待推导就遗漏返回。下面把泛型返回值与分支、循环组合。

\par
\medskip
\begin{lstlisting}[style=chali]
T fromLocal<T>(T value)
{
    return
    {
        T local = value;
        yield local;
    };
}

T chooseValue<T>(bool first, T left, T right)
{
    return if (first)
    {
        T local = left;
        yield local;
    }
    else
    {
        yield right;
    };
}

void main()
{
    String word = "generic";
    String copied = fromLocal<String>(word);
    int64 number = chooseValue<int64>(false, 4, 9);
    println("{} {} {}", word, copied, number);
}
\end{lstlisting}
\par
\medskip

\par
\medskip
\begin{lstlisting}[style=out]
generic generic 9
\end{lstlisting}
\par
\medskip

\subsection*{代码详解}

fromLocal 中的 T 同时用于参数、局部变量和返回值，实例化为 String 后整条路径都要满足字符串的复制、转移与清理规则。块内 local 通过 yield 形成块结果，再由 return 交回；word 后面仍被打印，说明不能为了返回值而非法提前使它失效。

chooseValue<int64> 接收 false，因此不执行 left 所在分支，直接由 else 交出 right 的 9。最终 word 与 copied 都显示 generic，number 为 9。泛型复用的是程序结构，具体类型仍参与每一条分支的检查，不会让模板中的局部变量变成无类型存储。

\section{实例化、代码体积与验证}

具体实例给后端提供类型信息，但更多实例也可能增加编译工作和产物体积。根据业务使用的组合验证行为，不把某一类型的结果当成所有基础类型、拥有值、引用和资源类型的证据。宏的类型参数与 \#res 推导见第 19 章。

\subsection*{小结}

类型参数决定使用什么类型，常量参数可决定数组形状，约束规定允许的操作，关联类型把实现选择传递给结果。它们提供静态关系，不是四种可互换的运行期数据。

泛型实例仍遵守复制、移动、借用和清理规则。遇到模板中的错误，先将本次实参代入，再像普通程序一样检查具体路径。下一章把变化的“行为”也作为函数值传递。

\section{习题}

\begin{enumerate}

\item 指出 Buffer 的两个参数分别控制什么。

\item 约束与强制转换有何区别？

\item 给泛型返回程序增加一个分支。

\item 选择三种元素类型验证泛型容器。

\end{enumerate}

\subsection*{参考解答与验收要点}

\begin{enumerate}

\item T 控制元素类型，N 控制常量维度。

\item 约束规定可用能力，转换不能替代能力实现。

\item 确保新增路径也返回声明结果类型。

\item 例如整数、String、带资源的结构体，分别核对清理与复制规则。

\end{enumerate}

下一章：\hyperref[ch:15]{第 15 章}。

\section{番外：郑人买泛型}

改写自《韩非子·外储说左上》郑人买履一则。

郑人有欲买履者，先自度其足，而置之其坐。及至市，乃忘操之。见一泛型鞋铺，谓店主曰：“吾尺在家，汝只须照吾脚做，举世之脚亦照吾脚做。”

店主曰：“函数可通用，尺寸须传入。君欲买的是全家鞋，何独固定一家之长？”郑人曰：“宁信吾已验证之具体类型，无自信于参数。”

于是给小儿做巨履，给老母做紧靴，给邻犬也套上一双。郑人解释：“接口统一，部署方便。”邻犬一口咬住了发布说明。

店主劝他留下制作的方法，将尺寸和材质作为参数；需要比较，便写比较的约束，需要缝合，便写缝合的约束。郑人大喜，捧来一块石头：“既称通用，此物也可制履乎？”

店主指着错误信息曰：“它没有你要求的能力。”郑人不悦，谓泛型无用，遂反归取度。及反，市罢，编译亦罢。

人曰：“何不试之以足？”曰：“宁信类型实参，不敢忘记约束。”店主听了，觉得虽走了远路，这回倒终于学会了半句正确的话。

\subsection*{技术对照}

泛型复用算法结构，并通过类型参数和约束表达所需能力。推导失败需要明确类型信息或调整调用，不等于运行期尝试任意类型。
